\documentclass{article}
\linespread{1.3}
\usepackage[margin=50pt]{geometry}
\usepackage{amsmath, amsthm, amssymb, amsthm, tikz, fancyhdr, float, graphicx, spverbatim, systeme}
\pagestyle{fancy}
\renewcommand{\headrulewidth}{0pt}
\newcommand{\changefont}{\fontsize{15}{15}\selectfont}

\fancypagestyle{firstpageheader}{
  \fancyhead[R]{
    \changefont
    \parbox[t]{4cm}{ % Adjust width as needed
      Michael Huang\\
      EN.625.633.81\\
      Homework N
    }
  }
}

\begin{document}

\thispagestyle{firstpageheader}
{\Large 

\section*{1.}

\subsection*{(a)} 

We are given \(X\) and \(Y\) with joint probability density function \(f(x, y) = 2 \exp\{-(x + 2y)\}\), given \(0 < x < \infty, 0 < y < \infty\). \\ \\
To find \(P(X < Y)\), we can take the integral of the pdf over the range where \(x < y\), so we would integrate from \(0 < x < y\):
\[
  P(X < Y) = \int_{y=0}^{\infty} \int_{x=0}^{y} f(x, y) dx dy
\]
\[
  = \int_{y=0}^{\infty} \int_{x=0}^{y} 2 \exp\{-(x + 2y)\} dx dy
\]
\[
  = \int_{y=0}^{\infty} \int_{x=0}^{y} 2 e^{-x} \cdot  e^{-2y} dx dy
\]
\[
  = \int_{y=0}^{\infty} 2e^{-2y} \int_{x=0}^{y} e^{-x} dx dy
\]
\[
  = \int_{y=0}^{\infty} 2e^{-2y} \cdot -e^{-x}|_{x=0}^{y} dy
\]
\[
  = \int_{y=0}^{\infty} 2e^{-2y} \cdot (-e^{-y} - (-e^{0})) dy
\]
\[
  = \int_{y=0}^{\infty} 2e^{-2y} \cdot (-e^{-y} + 1) dy
\]
\[
  = \int_{y=0}^{\infty} 2e^{-2y} - 2e^{-3y} dy
\]
\[
  = 2 \cdot \frac{e^{-2y}}{-2} - 2 \cdot \frac{e^{-3y}}{-3} |_{y=0}^{\infty}
\]
\[
  = -e^{-2y} - \frac{2e^{-3y}}{-3} |_{y=0}^{\infty}
\]
\[
  = 0 - (-e^0) - (0 - \frac{2e^{0}}{-3})
\]
\[
  P(X < Y) = 0 + 1 - 0 - \frac{2}{3} = 1 - \frac{2}{3} = \boxed{\frac{1}{3}}
\]

\subsection*{(b)}

To find \(P(X > \frac{1}{2})\), we need to first find the marginal of \(f(x, y)\) by integrating out \(y\):
\[
  f(x) = \int_{y=0}^{\infty} f(x, y) dy
\]
\[
  = \int_{y=0}^{\infty} 2 \exp\{-(x + 2y)\} dy = \int_{y=0}^{\infty} 2 e^{(-x - 2y)} dy
\]
\[
  = 2e^{-x} \int_{y=0}^{\infty} e^{-2y} dy
\]
\[
  = 2e^{-x} \cdot -\frac{e^{-2y}}{2} |_{y=0}^{\infty}
\]
\[
  = 2e^{-x} \cdot (-\frac{e^{-\infty}}{2}-(-\frac{e^0}{2}))
\]
\[
  = 2e^{-x} \cdot (0 + \frac{1}{2})
\]
\[
  f(x) = e^{-x}
\]

Now, we can properly integrate the marginal pdf over \(\frac{1}{2} < x < \infty\)
\[
  P(X > \frac{1}{2}) = \int_{x=\frac{1}{2}}^{\infty} f(x) dx
\]
\[
  = \int_{x=\frac{1}{2}}^{\infty} e^{-x} dx
\]
\[
  = -e^{-x} |_{x=\frac{1}{2}}^{\infty}
\]
\[
  = -e^{-\infty} - (-e^{-\frac{1}{2}}) 
\]
\[
  P(X > \frac{1}{2}) = 0 + e^{-\frac{1}{2}} = \boxed{e^{-\frac{1}{2}}}
\]



}
{\Large 

\section*{2.}

\subsection*{(a)} 

We are given \(f_X(x) = \frac{3}{16}x^2\) with \(-2 < x < 2\), and \(S = X_1 + X_2 + \dots + X_{24} + X_{25}\), with \(X_i\) all being independent with density \(f\). We can try to estimate using the Central Limit Theorem with the \(i\) independent random variables composing \(S\). To apply this, we would apply \(Z = \frac{S - E[S]}{\text{Std}(S)}\), so let's do some calculations for expectation and variance accordingly:
\[
  E[X] = \int_{-2}^{2} x \cdot f_X(x) dx
\]
\[
  = \int_{-2}^{2} x \cdot \frac{3}{16}x^2 dx
\]
\[
  = \int_{-2}^{2} \frac{3}{16}x^3 dx
\]
\[
  = \frac{3}{16 \cdot 4}x^4 |_{-2}^{2}
\]
\[
  = \frac{3 \cdot 2^4}{64} - \frac{3 \cdot -2^4}{64}
\]
\[
  E[X] = 0
\]
\[
  E[X^2] = \int_{-2}^{2} x^2 \cdot f_X(x) dx
\]
\[
  = \int_{-2}^{2} x^2 \cdot \frac{3}{16}x^2 dx
\]
\[
  = \int_{-2}^{2} \frac{3}{16}x^4 dx
\]
\[
  = \frac{3x^5}{16 \cdot 5} |_{-2}^{2}
\]
\[
  = \frac{3 \cdot 2^1}{5} - \frac{3 \cdot -2^1}{5}
\]
\[
  = \frac{6}{5} - \frac{-6}{5} 
\]
\[
  E[X^2] = \frac{12}{5}
\]
\[
  \text{Var}(X) = E[X^2] - (E[X])^2
\]
\[
  = \frac{12}{5} - 0^2
\]
\[
  \text{Var}(X) = \frac{12}{5}
\]
So for \(S\), we can accordingly say that 
\[
  E[S] = 25 \cdot 0 = 0
\]
\[
  \text{Std}(S) = \sqrt{25 \cdot \frac{12}{5}} = \sqrt{60} = 2\sqrt{15}
\]
To find \(P(\frac{1}{25} \geq .05) = P(S \geq 1.25)\): 
\[
  P(Z \geq \frac{S - E[S]}{\text{Std}(S)})
\]
\[
  = P(Z \geq \frac{1.25 - 0}{2\sqrt{15}})
\]
\[
  = P(Z \geq \frac{1.25}{2\sqrt{15}}) = P(Z \geq 0.16137430609)
\]
\[
  = 1 - \Phi(0.16137430609)
\]
\[
  P(\frac{1}{25} \geq .05) \approx 1 - 0.56356 = \boxed{0.43644}
\]

\subsection*{(b)}

To evaluate the quality of the approximation, let's first evaluate the actual variables. \(f_X(x)\) is symmetric about 0, with the expectation being exactly that, as \(S\) is the sum of all these \(X_i\). The approximation is about the mean / average \(X_i\), with \(\frac{S}{25}\) being this with \(n = 25\). So we are determining the probability of the mean \(X_i\) being greater than 0.05. The range is from \([-2, 2]\), with 0.05 being a relatively small value and the distribution having lighter distribution toward the edges of the parabola distribution, with the probability being slightly lower than 0.5 makes sense, as we are only slightly in the positive side of the distribution. So the approximation seems to generally be in the right area, and isn't egregiously bad by any means.

}
{\Large 

\section*{3.}

\subsection*{(a)} 



\subsection*{(b)}



\subsection*{(c)}



}
{\Large 

\section*{4.}

\subsection*{(a)} 

We are given a Poisson process with \(\lambda = 10\), and that two customers arrived in the first 5 minutes. As the arrival times are uniformly distributed, we have \(N(t) = n \) in time \([0, s]\) follows the Binomial with \(p = \frac{s}{t} = \frac{2}{5}\). We can easily calculate to calculate the probability that both arrived in the first 2 minutes as 
\[
  P(N(s) = l | N(t) = k) = \binom{k}{l}(\frac{s}{t})^l(1 - \frac{s}{t})^{k-l}
\]
\[
  P(N(s) = 2 | N(5) = 2) = \binom{2}{2}(\frac{2}{5})^2(1 - \frac{2}{5})^{2-2}
\]
\[
  = 1 \cdot \frac{4}{25} \cdot \frac{3}{5}^0
\]
\[
  P(N(s) = 2 | N(5) = 2) = \boxed{\frac{4}{25}}
\]

\subsection*{(b)}

To determine the probability that at least one arrived in the first two minutes, we can determine this by first determining the probability that no customers arrived in the first two minutes:
\[
  P(N(s) = l | N(t) = k) = \binom{k}{l}(\frac{s}{t})^l(1 - \frac{s}{t})^{k-l}
\]
\[
  P(N(2) = 0 | N(5) = 2) = \binom{2}{0}(\frac{2}{5})^0(1 - \frac{2}{5})^{2-0}
\]
\[
  = 1 \cdot 1 \cdot \frac{3}{5}^2
\]
\[
  P(N(2) = 0 | N(5) = 2) = \frac{9}{25}
\]
Now that we've found the probability that no customers arrived in the first two minutes, we can take the complement to determine the probability of having at least one customer arriving:
\[
  1 - \frac{9}{25} = \boxed{\frac{16}{25}}
\]

}
{\Large 

\section*{5.}

\subsection*{(a)} 

Let's set up the transition matrix accordingly:
\[
  P = 
  \begin{pmatrix}
    0.7 & 0.3 \\
    0.2 & 0.8 
  \end{pmatrix}
\]
and solve for the limiting distribution by solving \(\pi = \pi P\):
\begin{align*}
  \pi_A &= 0.7\pi_A + 0.2\pi_B \\
  \pi_B &= 0.3\pi_A + 0.8\pi_B
\end{align*}
\begin{align*}
  0.3\pi_A &= 0.2\pi_B \\
  \pi_B &= 0.3\pi_A + 0.8\pi_B
\end{align*}
\begin{align*}
  1.5\pi_A &= \pi_B \\
  \pi_B &= 0.3\pi_A + 0.8\pi_B
\end{align*}
\begin{align*}
  1.5\pi_A &= \pi_B
\end{align*}
\begin{align*}
  1.5\pi_A &= \pi_B \\
  \pi_A + \pi_B &= 1
\end{align*}
\[
  \pi_A + 1.5\pi_A = 1
\]
\[
  2.5\pi_A = 1
\]
\[
  \pi_A = \frac{2}{5}
\]
\[
  \pi_B = 1 - \frac{2}{5} = \frac{3}{5}
\]
So we have the limiting market shares as \(\boxed{A = 0.4, B = 0.6}\)

\subsection*{(b)}

We perform the same process with the new matrix: 
\[
  P = 
  \begin{pmatrix}
    0.7 & 0.15 & 0.15 \\
    0.1 & 0.8 & 0.1 \\
    0.05 & 0.05 & 0.9 \\
  \end{pmatrix}
\]
And now solving \(\pi = \pi P\):
\begin{align*}
  \pi_A &= 0.7\pi_A + 0.1\pi_B + 0.05\pi_C \\
  \pi_B &= 0.15\pi_A + 0.8\pi_B + 0.05\pi_C \\
  \pi_C &= 0.15\pi_A + 0.1\pi_B + 0.9\pi_C
\end{align*}
\[
  \pi_C - \pi_B = 0.15\pi_A + 0.1\pi_B + 0.9\pi_C - (0.15\pi_A + 0.8\pi_B + 0.05\pi_C)
\]
\[
  \pi_C - \pi_B = -0.7\pi_B + 0.85\pi_C
\]
\[
  0.15\pi_C = 0.3\pi_B
\]
\[
  \pi_C = 2\pi_B
\]
\[
  \pi_A - \pi_B = 0.7\pi_A + 0.1\pi_B + 0.05\pi_C - (0.15\pi_A + 0.8\pi_B + 0.05\pi_C)
\]
\[
  \pi_A - \pi_B = 0.55\pi_A - 0.7\pi_B 
\]
\[
  0.45\pi_A = 0.3\pi_B
\]
\[
  3\pi_A = 2\pi_B
\]
\[
  \pi_A + \pi_B + \pi_C = 1
\]
\[
  \frac{2}{3}\pi_B + \pi_B + 2\pi_B = 1
\]
\[
  \frac{11}{3}\pi_B = 1
\]
\[
  \pi_B = \frac{3}{11}
\]
\[
  \pi_C = 2\pi_B = 2 \cdot \frac{3}{11} = \frac{6}{11}
\]
\[
  3\pi_A = 2 \cdot \frac{3}{11}
\]
\[
  \pi_A = 2 \cdot \frac{1}{11} = \frac{2}{11}
\]
So the limiting market shares are \(\boxed{A = \frac{2}{11}, B = \frac{3}{11}, C = \frac{6}{11}}\)

}
{\Large 

\section*{6.}

\subsection*{(a)} 



\subsection*{(b)}



}
{\Large 

\section*{7.}



}
{\Large 

\section*{8.}



}
{\Large 

\section*{9.}

We have \(m\) cumulative distribution functions \(F_1(x), F_2(x), \dots, F_m(x)\), with each \(F_j(x)\) easily generatable. We aim to generate a random variable with the cdf \(F(x) = \sum_{j=1}^{m} p_j F_j(x)\) where \(\sum_{j=1}^{m} p_j = 1\). Essentially, we have a function \(F(x)\) which consists of finding of the weighted probability of each of the \(F_j(x)\) values. Intuitively, we need to distribute the values according to the probabilities \(p_j\), which we know they represent accordingly as they sum up to 1. We can easily see that this matches up with the law of total probability, i.e. that 
\[
  P(X \leq x) = \sum_{j=1}^{m} P(X \leq x | j=i) \cdot P(j=i)
\]
This is essentially equivalent to an event contingent on selecting some \(j = i\). So naturally, we would want to find a way to randomly generate the variable based on \(j\) and find \(p_j\) accordingly in order to find the corresponding \(F_j(x)\). Conveniently, an intuitive way to distribute this is to, as we've learned, randomly generate a value from the Uniform function on range \([0, 1]\), and then determine the corresponding \(j\) value based on the probability \(p_j\) accordingly, which we'll divide the range into accordingly, and then easily generate the corresponding value \(F_j(x)\). So putting this together, we'd do the following: \\
- Divide up the range \([0, 1]\) into different segments based on the \(p\)-values \([0, p_1, p_1 + p_2, p_1 + p_2 + p_3, \dots, \sum_{j=1}^m p_j]\) \\ 
- Generate value \(U \sim \text{Unif}(0, 1)\) \\
- Determine with where \(U\) falls on the range of values to determine the corresponding \(p_j\) value, where \(j\) represents the last \(p\)-value added to the upper bound of the \(p\)-value ranges. \\
- Using the corresponding \(j\), then generate the random variable \(X\) according to \(F_j(x)\). For this, we can then use the inverse-transform method, specifically, first generate some other \(U' \sim \text{Unif}(0, 1)\) and use that \(U'\) to determine the final \(X = F_j^{-1} (U')\).

}

\end{document}